\section{Local jet extraction and exact masked splitting}
\label{sec:local-splitting}
A transferred local field is not yet a coefficient field. Its low-order
jet has to be extracted, and subtracting that jet changes the support of
its healthy-set indicator. The difference of indicators is retained as
an explicit flagged response. We give the extraction and this exact
splitting before applying it to the one-step integral.

Fix a finite lattice, an anchor $o$, a degree $d$, and a coefficient
frame $V_d$ for the gauge-invariant curvature-jet polynomials of weighted
degree at most $d$. A curvature derivative of order $j$ has weight
$j+2$. Write $\mu(v)=\sum_\alpha|v_\alpha|$, $e_0$ for the unit
coordinate, $v^\circ=v-v_0e_0$, and $D_t$ for the diagonal map multiplying
a coordinate with $r$ curvature factors by $t^r$.
The budget consists of a threshold $0<t\le1$, derivative order $q\ge d$,
spatial power $p\ge d$, positive support exponent $u$, and integer
covering radius $B\ge1$.

A regular field on a support $X$ is a measurable, gauge-invariant
function depending only on its links, whose exponential-coordinate germ
is $C^q$ at every configuration healthy on $X$. Its ordered left
 derivatives exist through order $q$, including at the threshold boundary.
Write $\|f\|_{o,X}\le c$ when $c\ge0$ and all these derivatives of
orders at most $q$ have modulus at most $c$, for direction lists whose
every entry satisfies
\[
       |A(e)|\le t(1+\operatorname{dist}(e_-,o))^p.
\]
In particular the zeroth-order estimate bounds $f$ on the healthy set.
No derivatives of a healthy-set indicator occur in this definition.

\subsection{Reading the unmasked jet}
Let $H_o(J,h)$ be the local test holonomy for a curvature jet $J$,
and write $h\cdot J$ for weighted scaling, with component
$h^{j+2}J_j$. Its defining local identities are
\[
 H_o(J,0)=1,\qquad H_o(J,h)=H_o(h\cdot J,1),\qquad
 H_o(gJ,h)=g\,H_o(J,h)\,g^{-1}.
\]
Here only the components of weights at most $d$ are retained on the
right. The finite-dimensional map $Z\mapsto H_o(Z,1)$ is analytic near
$Z=0$ in the genuine group exponential chart. These properties of the
chosen test holonomy, and the existence of the local polynomial section
below, are the geometric inputs used in this section.

Here a support $X$ is nonwrapping at $o$ precisely when
$10\operatorname{dist}(x,o)<n$ for every $x\in X$, where $n$ is
the lattice side. For such a support, define $\Pi_{o,X}f$ by requiring that
for every $J$, the functions
\[
       h\longmapsto f(H_o(J,h)),\qquad
       h\longmapsto (\Pi_{o,X}f)(h\cdot J)
\]
have the same derivatives at zero through order $d$.
For a wrapping support the convention is
$\Pi_{o,X}f=f(1)\,1$.

We first prove existence in the nonwrapping case. In the finite space
of truncated jets put $F(Z)=f(H_o(Z,1))$. The identity configuration
is healthy, and the local logarithm of $H_o(Z,1)$ is analytic and
vanishes at zero. Thus $F$ is $C^d$ near zero. Its Taylor polynomial
\[
        T(Z)=\sum_{r=0}^d\frac1{r!}D^rF(0)[Z,\ldots,Z]
\]
is a polynomial in the real coordinates of the curvature jets.
Delete its monomials of weighted degree greater than $d$, obtaining
$P$. The ordinary Taylor polynomial has the same derivatives as $F$
through order $d$. Composing with the polynomial curve $h\cdot J$
preserves that equality: in the chain rule for a derivative of order
$k\le d$, only derivatives of $F$ of orders at most $k$ enter.
A deleted monomial becomes a multiple of $h^r$ with $r>d$, so contributes
nothing through order $d$. Hence $P$ has exactly the required test jets.

Gauge invariance of $f$ and covariance of $H$ give the same test jets
for $P(gJ)$ and $P(J)$. Their difference has weighted degree at most
$d$. For each fixed $J$, all its coefficients as a polynomial in $h$
vanish, since its derivatives at zero vanish. Taking $h=1$ proves
$P(gJ)=P(J)$. Thus $P$ belongs to the gauge-invariant label space.
The same argument applied to two candidates proves uniqueness.
Subtraction and scalar multiplication of the test identities now show
that $\Pi_{o,X}$ is linear on regular fields with a fixed support.
The wrapping convention has these properties as well.

For later use set
\[
                  v(f)=D_t\operatorname{coord}(\Pi_{o,X}f).
\]
The uniform estimate for this map follows from a precise local bound
on test holonomies. Namely, for every fixed $J,d$ and $r\le d$, suppose
the logarithmic test curve $Z(h)=\log H_o(tJ,h)$ satisfies
\begin{equation}\label{eq:local-test-log-input}
 |Z^{(r)}(0)(e)|\le a_r(J,d)\,t
                    (1+\operatorname{dist}(e_-,o))^{\max(r-1,0)}.
\end{equation}
The constants are independent of the period, anchor, support and
threshold. We prove this estimate for the actual radial test below;
first we explain precisely how it is used.

Here is why it gives a constant $R_d\ge1$ such that
\begin{equation}\label{eq:local-read-price}
                         \mu(v(f))\le R_d c.
\end{equation}
First bound a mixed exponential-coordinate derivative of order $k$
on directions $w_1,\ldots,w_k$ bounded by fixed multiples
$a_1,\ldots,a_k$ of the allowed direction weight. Restrict the symmetric
$k$th Fr\'echet derivative at zero to their span, using the map
$Lx=\sum_i x_iw_i$ from $\mathbb R^k$. Its diagonal at $x$ is bounded
by $(1+\sum_i|x_i|a_i)^k c$: normalize $Lx$ by this positive factor
and use the repeated same-direction left derivative along
$\exp(sLx)$. Such a repeated derivative equals the derivative of
$f(\exp(sLx))$; the one-parameter group identity proves this by
induction. Only the Fr\'echet derivative in this finite-dimensional
chart is being treated as symmetric.

Diagonal evaluations separate symmetric tensors on $\mathbb R^k$.
A finite number of them therefore bound any fixed mixed evaluation.
For completeness, if separating linear functionals $\ell_i$ on a
finite-dimensional space obey $|\ell_i A|\le b_i c$, but no uniform
norm bound follows, choose $A_n,c_n$ with
$\|A_n\|>(n+1)c_n$. A subsequence of $A_n/\|A_n\|$ converges on the
unit sphere. Every $\ell_i$ vanishes on its limit, a contradiction.
Apply this argument to the symmetric tensors just described. The
resulting mixed bound depends on $k$ and the $a_i$, not on the lattice.

The finite-order chain rule for $f(\exp Z(h))$, together with
\eqref{eq:local-test-log-input} and $p\ge d$, now bounds each fixed
jet probe by $C_{J,k}c$. These probes, for $k\le d$, separate $V_d$:
a weighted polynomial all of whose test derivatives vanish is zero,
by the uniqueness argument above. Apply the same finite-dimensional
argument to $V_d$ to obtain \eqref{eq:local-read-price}. The weights
are exactly $D_t$, since evaluating a polynomial at $tJ$ multiplies
a monomial with $r$ curvature factors by $t^r$.
For a wrapping support only $f(1)$ remains, and $|f(1)|\le c$.

\subsection{Uniform bounds for the logarithmic test links}
We establish \eqref{eq:local-test-log-input} by the same rescaling
used to construct the test. A truncated radial connection is a finite
sum of homogeneous connection tensors $J_j$ of degree $j+1$.
On a positive link starting at a covering-space point $x$ in direction
$e_i$, its coefficient at test scale $h$ and amplitude $t$ is
\[
       \sum_{j=0}^{d-2}t h^{j+2}
                    J_j(e_i)[x+s e_i,\ldots,x+s e_i],
                  \qquad 0\le s\le1.
\]
An empty sum is allowed. The link is the endpoint of the
right-multiplication equation $U'=UA$, $U(0)=1$, with this
Lie-algebra-valued coefficient. The local logarithm is taken only
near $h=0$.

We need an elementary parameter regularity fact for this equation.
On continuous matrix paths with the supremum norm define
$T_A V(s)=\int_0^s V(r)A(r)\,dr$. Its norm is at most
$\|A\|_\infty$. If $A$ depends linearly on a finite-dimensional
parameter $z$, then $T_A$ is a continuous linear function of $z$.
For $z$ near zero, $\|T_A\|<1$, and the integral equation has the
unique solution
\[
                         U=(I-T_A)^{-1}1
                           =\sum_{k=0}^\infty T_A^k1.
\]
The convergent operator power series makes this an analytic function
of $z$. The given link solution satisfies the same integral equation,
so agrees with it. Endpoint evaluation and the analytic local group
logarithm therefore give an analytic logarithmic endpoint map $E(z)$
near zero, with $E(0)=0$. This argument applies to the finite-dimensional
space of all polynomial-in-$s$ connection profiles of the indicated
degrees, not only to profiles coming from a particular link.

Put $R=1+\operatorname{dist}(e_-,o)$, $r=R^{-1}$ and
$y=r\,\operatorname{lift}(e_--o)$. The centered lift has each coordinate
at most $R$ in modulus, so $|y|_2\le2$, while $0<r\le1$.
Consider the polynomial profile
\[
       P(y,r,a,h)(s)=\sum_{j=0}^{d-2}a h^{j+2}
                    J_j(e_i)[y+rse_i,\ldots,y+rse_i].
\]
For all parameters $(y,r,a)$ the profile is zero at $h=0$.
Consequently
\[
       Q_k(y,r,a)=\left.\frac{d^k}{dh^k}
                                E(P(y,r,a,h))\right|_{h=0}
\]
is smooth in $(y,r,a)$: near each point it is obtained by
composing a polynomial map with the analytic germ $E$, then
differentiating a fixed finite number of times. Also $Q_k(y,r,0)=0$.
On the compact set $|y|_2\le2$, $0\le r,a\le1$, its derivative
with respect to $a$ is bounded. Integrating that derivative from
zero to $a$ gives $|Q_k(y,r,a)|\le C_{J,d,k,i}a$.
This compactness is in the fixed profile-parameter space and has
no dimension dependence on the lattice.

Homogeneity of $J_j$ shows that the actual link at scale $hr$ and
amplitude $t$ has profile $P(y,r,tr,h)$: both coefficients are
$t h^{j+2}r^{j+2}J_j(e_i)[x+se_i,\ldots,x+se_i]$.
Uniqueness of the link equation identifies their endpoints.
Differentiating this identity at zero gives
\[
       r^k|Z^{(k)}(0)(e)|\le C_{J,d,k,i}tr.
\]
For $k\ge1$ this is precisely
$|Z^{(k)}(0)(e)|\le C_{J,d,k,i}t R^{k-1}$.
For $k=0$ the logarithm is zero. Taking the maximum over the four
directions proves \eqref{eq:local-test-log-input}. The argument
uses only derivatives at zero; it asserts no common analytic radius
for links arbitrarily far from the anchor.

\subsection{The support correction}
Choose a linear local section of the labels, supported on a fixed
stencil $Y$ containing $o$. Its placed field $P_t(v)$ represents
$D_t^{-1}v$, has the prescribed jet, and satisfies
\begin{equation}\label{eq:local-section-price}
              \|P_t(v)\|_{o,Y}\le w\mu(v),\qquad w\ge0.
\end{equation}
The stencil is assumed to fit without wrapping. The section sends
$e_0$ to the constant one. These are hypotheses about the chosen
section; they do not follow from the abstract existence of $\Pi$.
Let $\chi_X$ be the indicator that $X$ is healthy, and put
\[
       p_t(v)=v_0+\chi_Y P_t(v^\circ),\qquad Z=X\cup Y.
\]
Then $|p_t(v)|\le(1+w)\mu(v)$ everywhere, and
$p_t(v)=P_t(v)$ when $Y$ is healthy. Define
\begin{align}
       g&=f-P_t(v(f)),\\
       r&=\chi_X f-p_t(v(f))-\chi_Zg.
       \label{eq:local-rebate}
\end{align}
These definitions give the exact, pointwise identity
\begin{equation}\label{eq:local-atom-identity}
                  \chi_X f=p_t(v(f))+\chi_Zg+r.
\end{equation}
The field $g$ is regular on $Z$. It has zero identity value and,
when $Z$ is nonwrapping, zero test jets through order $d$.
Indeed $Y$ is nonwrapping, so $Z$ is nonwrapping exactly when $X$
is; subtract the two equal jets in that case. In the wrapping case
$v(f)=f(1)e_0$, and the same subtraction gives the zero identity value.

The response $r$ is measurable, gauge invariant, and determined by
the links and final flags in $Z$. It vanishes whenever $Z$ is healthy,
by \eqref{eq:local-atom-identity} and $p_t(v)=P_t(v)$ there.
Put $A_d=1+(1+w)R_d$. Then
\begin{equation}\label{eq:local-atom-cost}
             \|g\|_{o,Z}\le A_dc,\qquad |r(U)|\le A_dc.
\end{equation}
The first estimate is the triangle inequality using
\eqref{eq:local-read-price}--\eqref{eq:local-section-price}.
For the second, $r=0$ on the healthy set of $Z$. On its complement
$\chi_Z=0$, so
$|r|\le|\chi_X f|+|p_t(v(f))|\le c+(1+w)R_dc$.
This proves the global bound without differentiating any mask.

\subsection{Putting the response into the existing gas}
Let $K_j$ be the original flagged activities, with supports $P_j$
and nonempty flag inventories $I_j\subseteq P_j$. An admissible
finite family has disjoint supports and its inventories cover
exactly the final flags. Let $\mathcal G(U)$ be its gas sum.
For each response $r_i$ supported on $Z_i$ and each nonempty
pairwise disjoint family $q$ of old activities meeting $Z_i$, assume
a distinct available address with support and inventory
\[
       Z_i\cup\bigcup_{j\in q}P_j,\qquad \bigcup_{j\in q}I_j.
\]
The address's additional eligibility predicate must hold whenever
its support has exactly this flag inventory. No zero value of its
old activity is required. Insert at this address the coefficient
\[
       r_i(U)\prod_{j\in q}K_j(U)
\]
when eligible, and zero otherwise. Denote the resulting source
array by $\mathcal K(r)$. This is a linear operation in $r$, with
addresses fixed before their values are evaluated.

Its exact marked-gas sum is
\begin{equation}\label{eq:local-gas-recovery}
       \mathcal G^{[1]}(\mathcal K(r))(U)
                     =\Bigl(\sum_i r_i(U)\Bigr)\mathcal G(U).
\end{equation}
To prove this, take a nonzero term on the right, consisting of
$i$ and an admissible old family $b$. Because $r_i$ vanishes on
healthy $Z_i$, some final flag lies in $Z_i$. Exact inventory
coverage implies that the subfamily
$q=\{j\in b:P_j\cap Z_i\ne\varnothing\}$ is nonempty.
Replace precisely this subfamily by the designated new address.
Every remaining support is disjoint from the new support;
the inventories still cover the same flags. This produces an
admissible pointed family on the left, with identical value.
Conversely, expand the distinguished new address into its encoded
family $q$. The other supports are disjoint from its total support,
so $q$ is exactly the subfamily meeting $Z_i$. This recovers $i,b$
uniquely. Thus there is a value-preserving bijection of nonzero
terms. Absolute summability, established next, permits reindexing.

Let $C_B(X)$ be the covering number used in the local norm and set
\[
       m(X,a)=e^{uC_B(X)}a,\qquad
       m_+(X,a)=e^{(u+2)C_B(X)}a.
\]
Suppose $|K_j|\le b_j$, $b_j\ge0$, and, for a site price $\varsigma\ge0$
(not to be confused with the coefficient vectors $v$ above),
$\sum_{j:x\in P_j}m(P_j,b_j)\le\varsigma$ at every site $x$.
For one response with $|r_i|\le a_i$, the inserted array has total
mass at most
\begin{equation}\label{eq:local-merge-cost}
           m(Z_i,a_i)\bigl(e^{|Z_i|\varsigma}-1\bigr).
\end{equation}
Indeed the total mass of old activities meeting $Z_i$ is at most
$|Z_i|\varsigma$, by summing their site bounds over $Z_i$.
The mass of the union is at most the product of the masses, since
each old support meets $Z_i$ and their covering numbers are
subadditive along such unions. Summing products over nonempty
finite subsets is bounded by $e^{\sum b}-1$; dropping the
disjointness restriction only increases this nonnegative sum.
This proves \eqref{eq:local-merge-cost}. If the sum of the right sides of \eqref{eq:local-merge-cost}
over $i$ is finite, summing these bounds proves absolute domination
of the inserted array and hence of its pointed gas series.
The reserved-mass hypothesis below will establish this finiteness. The old gas series is absolutely convergent
under the same site bound on the finite lattice. This justifies
\eqref{eq:local-gas-recovery}.

\subsection{The two-unit reserve in the support exponent}
Let $T$ bound $C_B(Y)$ and put $h_B=(2B+1)^4$.
The elementary covering estimates used here are
\[
 C_B(X\cup Y)\le C_B(X)+T\quad(o\in X\cap Y),
       \qquad |X|\le h_B C_B(X).
\]
Set $F_d=e^{(u+2)T}A_d$. Then
\begin{equation}\label{eq:local-join-price}
               m_+(X\cup Y,A_dc)\le F_d m_+(X,c).
\end{equation}
This follows immediately by inserting the first covering inequality
in the exponential. If $h_B\varsigma\le1$, then for $a\ge0$,
\begin{equation}\label{eq:local-meeting-reserve}
       m(X,a)(e^{|X|\varsigma}-1)\le h_B\varsigma\,m_+(X,a).
\end{equation}
For $D=C_B(X)$, put $q_0=h_B\varsigma$. We have $|X|\varsigma\le Dq_0\le D$.
Using $e^x-1\le xe^x$ for $x\ge0$ and $D\le e^D$ gives
$e^{|X|\varsigma}-1\le q_0De^D\le q_0e^{2D}$, as required.

Consider now a raw expansion with seed $v_0$, regular fields $f_i$
on anchored supports $X_i$, and an existing forced source $k$.
Suppose $\|f_i\|_{o,X_i}\le c_i$ and
\[
       \sum_i m_+(X_i,c_i)\le a,\qquad
                      \operatorname{mass}(k)\le\kappa b,
       \qquad \kappa=\sqrt\varsigma>0,\quad h_B\kappa\le1.
\]
All constants $a,b$ are nonnegative, and $\varsigma$ is the site price
of the original flagged activities $K_j$ of the preceding subsection.
Form $v_i=v(f_i)$,
$g_i,r_i$ as above, and define the coefficient field by
\begin{equation}\label{eq:local-built-field}
       v=v_0+\sum_i v_i,\qquad
       \text{regular terms }g_i\text{ on }X_i\cup Y,\qquad
       \text{forced array }k+\mathcal K(r).
\end{equation}
The coordinate sum converges absolutely by
\eqref{eq:local-read-price} and $\sum_i c_i\le a$.
Its change is bounded by $R_da$. The regular remainder mass is at
most $F_da$. Furthermore $h_B\varsigma=h_B\kappa^2\le\kappa\le1$,
so \eqref{eq:local-meeting-reserve} and
\eqref{eq:local-join-price} bound the added forced mass by
$\kappa F_da$. Measuring forced mass in units of $\kappa$, so that
the forced array has normalized mass at most $b+F_da$, the normalized
total remainder bound is
\begin{equation}\label{eq:local-built-cost}
                              2F_da+b.
\end{equation}
The original density is preserved exactly: sum
\eqref{eq:local-atom-identity}, use linearity of $p_t$, and then
\eqref{eq:local-gas-recovery} for the response terms. Multiplication
by the common vacuum exponential preserves this identity.
Every sum is absolutely dominated by the bounds just proved.
In particular this construction uses the original vacuum activities
and the original threshold predicates throughout.

Finally, collect the responses with the same support:
$\psi(Z)=\sum_{i:X_i\cup Y=Z}r_i$. Their majorants obey
\[
       \sum_Z m_+\left(Z,\sum_{i:X_i\cup Y=Z}A_dc_i\right)
                                    \le F_da.
\]
Consequently $\psi$ is an anchored, gauge-invariant response vanishing
on healthy supports, with the stated reserved mass. Its insertion
profile is exactly the profile of $\mathcal K(r)$: grouping the
absolutely convergent terms by their encoded support and inventory
proves the equality. These are the response and profile identities
needed by the marked transport of Section~\ref{sec:local-transfer}.
The raw expansion and its bounds still have to be produced by the
one-step integral; this section proves the projection, mask correction,
and cost conversion once that expansion is given.
